\scanpage{22}
Que $\Omega$ soit Coh. Mac. de dim résulte des
formules précédentes et de la caractérisation des modules
de Coh. Mac. de dim $n$. Il s'ensuit
$H^i(\Omega)=0$ pour $i<n$, \deleted{\unsure{et}}\inserted{et pour $i>n$ d'après le th. 3.} D'après le th. 1,
le cor. 1 implique que $H^n(\Omega)$ est dualisant.

D'après le th. de dualité, on a
\[
H^i(\Omega)\xrightarrow{\sim}D(\operatorname{Ext}^{n-i}(\Omega,\Omega))
\]
-- qui montre que \deleted{$H^i(\Omega)=0$ si $i>n$.}
$\operatorname{Ext}^{n-i}(\Omega,\Omega)=0$ si $\deleted{\unsure{n-}}i\ne n$ i.e. $n-i\ne0$.
\deleted{\unsure{Cela montre aussi que}} Si $\deleted{\unsure{n-}}i=n$, on
trouve $D(\operatorname{Hom}(\Omega,\Omega))\simeq H^n(\Omega)$ qui est
dualisant i.e. qui est $\simeq D(A)$, d'où
$\operatorname{Hom}(\Omega,\Omega)\simeq A$. Ainsi \uline{tout endomorphisme
de $\Omega$ est donné par un scalaire}, déterminé de
\uline{façon unique}. \inserted{A fortiori\unsure{;|,}} $\Omega$ \uline{est \deleted{\unsure{donc}} un module fidèle}.

\uline{Corollaire 4} \inserted{$A,\Omega,n$ comme dessus.} $\Omega_{\unsure{A}}$ est de dim inj. $=n$, en
particulier on a
\[
\operatorname{Ext}^i(M,\Omega)=0\qquad\text{si}\quad i>n
\]
\uline{quel que soit} $M$

Il suffit de le vérifier pour $M$ de
type fini, et ça résulte alors du th. de dualité.

% margin note written sideways
\margin{\uline{Corollaire 5} Soit $A$ un anneau local de Coh. Mac. de dim $n$, $\Omega$ un module fondamental pour $A$, $M$ un $A$-module de type fini\unsure{.} Alors \deleted{\unsure{on a}} $\operatorname{cocodim}M\leq k\Longleftrightarrow\operatorname{Ext}_A^{\unsure{i}}(M,\Omega)=0$ pour $i>n-k$. Si $M$ est de dim Krull $k$, cela équivaut alors à : $M$ est Coh. Mac. et implique $\operatorname{Ext}_A^i(M,\Omega)=0$ pour $i\ne k$.}
% archive sheet number: 21

